\bumppoints{24}

\question
Naturally occurring fires cause occasional disturbance to an ecosystem. If humans stop these fires as soon as they are started, then

\begin{choices}
\choice species diversity will tend to remain unchanged, but the types of species present may change.
\choice species diversity in the ecosystem will tend to increase.
\correctchoice species diversity in the ecosystem will tend to decrease.
\choice the diversity of threatened and endangered species will be protected.
\end{choices}

\question
Which of the following is \emph{not} directly concerned with community ecology?

\begin{choices}
	\choice herbivory
	\correctchoice intraspecific competition
	\choice resource partitioning
	\choice predation
	\choice symbiosis
\end{choices}

\question
The dispersion pattern for individuals across the entire metapopulation shown in question \ref{ques:metapopulation} is most likely

\begin{choices}
	\choice even. 
	\correctchoice clumped. 
	\choice uniform. 
	\choice gene flow. 
	\choice random. 
\end{choices}

\question
The attribute that most influences population density and rates of immigration and emigration among habitat patches in a metapopulation are

\begin{choices}
	\choice patch size, habitat quality, and isolation of sink patches. 
	\choice patch size, patch isolation and species diversity. 
	\correctchoice patch size, patch isolation and habitat quality. 
	\choice habitat quality, size of sink patches and isolation of source patches. 
	\choice patch isolation, habitat quality, and size of source patches. 
\end{choices}

\question
In the experiment shown in the figure below, \textit{Balanus} was removed from the habitat shown on the left. Which of the following statements is a valid conclusion of this experiment?\\
\includegraphics[width=0.9\textwidth]{exam3_balanus}

\begin{choices}
	\choice If \textit{Chthamalus} were removed, \textit{Balanus}'s fundamental niche would become larger.
	\choice These two species of barnacle do not show competitive exclusion.
	\correctchoice The removal of \textit{Balanus} shows that the realized niche of \textit{Chthamalus} is smaller than its fundamental niche.
	\choice \textit{Balanus} can survive only in the lower intertidal zone because it is unable to resist dehydration.
	\choice \textit{Balanus} is inferior to \textit{Chthamalus} in competing for space on rocks lower in the intertidal zone. 
\end{choices}

\question
Most insects have two pairs of wings. Flies only have one pair of wings. The second pair of wings is reduced to tiny structures called halteres. The mutation(s) that led to the change from wings to halteres could have

\begin{choices}
\choice duplicated all of the \emph{Hox} genes in the fly genome. 
\choice resulted in paedomorphosis. 
\choice caused the loss of \textit{Antp} or \textit{Ubx}. 
\correctchoice  altered the expression of a \emph{Hox} gene. 
\choice caused expression of a pseudogene. 
\end{choices}

\question
Which of the following statements is consistent with the principle of competitive exclusion?  

\begin{choices}
\choice Two species with the same fundamental niche will exclude other competing species.   
\choice Evolution tends to increase competition between related species.   
\correctchoice Even a slight reproductive advantage will eventually lead to the elimination of the less well adapted of two competing species. 
\choice Bird species generally do not compete for nesting sites.  
\choice The density of one competing species will have a positive impact on the population growth of the other competing species. 
\end{choices}

\newpage

\question
Given the following values and assuming exponential population growth, what will be the total population size after one generation has elapsed? Round to nearest whole individual.\vspace{1ex}

\hspace{\leftmargin}%
\begin{minipage}{2in}
	$b = 0.09$ \qquad $N = 18,657$\\
	$d = 0.12$ \qquad $K = 15,000$\\
\end{minipage}

\begin{choices}
	\choice $560$ %dontshuffle 
	\choice $3,657$ %dontshuffle 
	\choice $15,000$ %dontshuffle 
	\correctchoice $18,097$ %dontshuffle 
	\choice $19,217$ %dontshuffle 
\end{choices}

\question
Which of the following is most likely to contribute to density-dependent regulation of populations?

\begin{choices}
\choice Earthquakes.
\choice Floods.
\choice Fires.
\correctchoice Intraspecific competition for nutrients.
\choice The removal of toxic waste by decomposers.
\end{choices}

\question
White-breasted nuthatches and brown creepers both eat insects that hide in the furrows of bark in hardwood trees. The brown creeper searches for insects by hunting from the bottom of the tree trunk toward the top, whereas the white-breasted nuthatch searches from the top of the trunk down. Their feeding areas rarely overlap on the tree trunk. These hunting behaviors best illustrate which of the following ecological concepts?

\begin{choices}
\choice interspecific competition
\correctchoice resource partitioning
\choice competitive exclusion
\choice keystone species
\choice intraspecific competition
\end{choices}

\question
Which block in the figure below is most consistent intermediate levels of disturbance in an ecosystem? (Different shades, including white, represent different species. The small squares indicate relative abundance of individuals present.)

\begin{multicols}{3}
	\begin{choices}
\choice \includegraphics{diversity_mcC} 
\choice \includegraphics{diversity_mcA} 
\correctchoice \includegraphics{diversity_mcE} 
\choice \includegraphics{diversity_mcD} 
\choice \includegraphics{diversity_mcB} 
\end{choices}
\end{multicols}

\question
A harmless fly looks like a dangerous wasp. The fly is using

\begin{choices}
	\choice warning coloration. 
	\choice cryptic coloration. 
	\choice camouflage. 
	\choice disruptive coloration. 
	\correctchoice mimicry. 
\end{choices}

